%==============================================================================
% COVER LETTER FOR JOURNAL OF PHYSICAL CHEMISTRY LETTERS
%==============================================================================

\documentclass[11pt,a4paper]{letter}
\usepackage[utf8]{inputenc}
\usepackage[T1]{fontenc}
\usepackage{amsmath,siunitx}

\address{
Department of Physics\\
Faculty of Science\\
University of Yaoundé I\\
Yaoundé, Cameroon
}

\signature{Steve Cabrel Teguia Kouam}

\begin{document}

\begin{letter}{
Dr. Anne B. McCoy\\
Editor-in-Chief\\
The Journal of Physical Chemistry Letters\\
American Chemical Society
}

\opening{Dear Dr. McCoy,}

We submit ``Quantum-Coherent Spectral Engineering: Non-Markovian Dynamics in Photosynthetic Energy Transfer under Engineered Photon Baths'' for consideration as a Letter in The Journal of Physical Chemistry Letters.

\textbf{What we found:} Spectral engineering of the incident photon bath extends excitonic coherence lifetimes by \SIrange{20}{50}{\percent} and enhances excitonic energy transfer efficiency by \SI{25}{\percent} through non-Markovian dynamics. Dual-band filtering at \SIlist{750;820}{\nano\meter} targets vibronic resonances in the Fenna--Matthews--Olson complex, producing polaron-dressed states with reduced effective reorganization energy and suppressed pure dephasing.

\textbf{Why it matters:} This work establishes \textit{spectral bath engineering} as a general principle for coherence control in open quantum systems. Unlike previous studies that treat environmental coupling as a constraint, we demonstrate that deliberate spectral shaping can \textit{enhance} rather than suppress quantum effects. The mechanism---vibronic-resonance-assisted transport---is applicable beyond photosynthesis to any light-driven quantum process.

\textbf{Why JPCL:} Our findings bridge quantum dynamics theory with potential experimental realization via programmable light sources or filtered illumination in 2DES experiments. The concise 5-page format is ideal for presenting this focused physical insight without extensive materials discussion.

\textbf{Key contributions:}
\begin{enumerate}
\item Demonstration of \SIrange{20}{50}{\percent} coherence lifetime extension via spectral filtering
\item Physical mechanism: vibronic resonance condition selectively suppresses pure dephasing
\item Coherence enhancement ($\eta = 0.20 \pm 0.04$) persists under physiological disorder ($\sigma = \SI{50}{\per\centi\meter}$) and temperature (\SIrange{285}{300}{\kelvin})
\item 12-test validation suite confirms computational accuracy (HEOM benchmark: \SI{1.8}{\percent} deviation)
\end{enumerate}

This manuscript is a computational/theoretical study. All simulations use the Process Tensor Hierarchy of Pure States (PT-HOPS) method with Spectrally Bundled Dissipators (SBD), a numerically exact framework for non-Markovian quantum dynamics. Experimental validation via 2DES under filtered illumination is a promising direction for future work.

All authors have approved this manuscript. We declare no conflicts of interest.

\textbf{Suggested reviewers:}
\begin{itemize}
\item Gregory Scholes (Princeton) -- gscholes@princeton.edu
\item Alán Aspuru-Guzik (Toronto) -- alan@aspuru.com
\item Martin Plenio (Ulm) -- martin.plenio@uni-ulm.de
\item Jianshu Cao (MIT) -- jianshu@mit.edu
\item Graham Fleming (Berkeley) -- gfleming@berkeley.edu
\end{itemize}

\textbf{Excluded reviewers:} None requested.

This manuscript is not under consideration elsewhere. Upon acceptance, we will transfer copyright to ACS.

Thank you for your consideration.

\closing{Sincerely,}

\fromsig{
Steve Cabrel Teguia Kouam\\
Department of Physics\\
University of Douala, Cameroon\\
Email: steve.teguia@facsciences-uy1.cm
}

\end{letter}
\end{document}
